\documentclass{article}
\title{Reliable Edge Monitoring under Partial Sensing}
\begin{document}
\maketitle
\begin{abstract}
Industrial condition monitoring often operates with missing or degraded sensors. We present a reliability-aware edge monitoring framework that supports partial sensing and selective maintenance decisions. On a held-out evaluation set, the method improves fault detection accuracy by 8.4\% while reducing expected calibration error to 0.031 and maintaining 82\% decision coverage. The model runs at 12 ms per sample on an edge device.
\end{abstract}

\section{Introduction}
Existing multimodal monitoring systems often assume complete sensing; however, field deployments face intermittent sensors, acquisition cost, and edge-compute constraints. Our contributions are: (1) a mask-conditioned representation for incomplete sensing, (2) calibrated reliability-aware fusion, and (3) deployment-oriented evaluation with selective maintenance.

\section{Method}
The proposed system encodes each available modality, conditions the representation on the observation mask, and fuses available evidence using learned reliability weights. A calibrated confidence head supports safe, maintain, and abstain decisions.

\section{Experiments and Results}
We compare against unimodal and multimodal baselines under complete and missing-sensor regimes. Ablation studies remove the mask adapter, reliability fusion, and calibration objective. Results are averaged over five runs and reported with mean +/- standard deviation. We report accuracy, expected calibration error (ECE), Brier score, selective risk, coverage, latency, and memory footprint. Source code and implementation details are available in the project repository.

\section{Limitations}
The current validation covers one industrial monitoring domain and does not establish universal performance under arbitrary domain shift. Failure cases include simultaneous severe degradation of all sensing channels.

\section{Conclusion}
The results support reliability-aware partial-sensing decisions under the tested deployment conditions while preserving edge feasibility.
\end{document}
